% ==============================================================================
% CAPÍTULO 9: PRUEBAS EN EL SOFTWARE (Orquestador del Capítulo)
% ==============================================================================
\preportadacapitulo{9}{Pruebas en el Software}

\newpage
\section{Capítulo 9: Pruebas en el Software}\label{cap:pruebas}

% [INSTRUCCIÓN]: Flujo de trabajo de Pruebas del PUDS. Contenido propio de esta versión.

\subsection{9.1. Pruebas}

% [PENDIENTE]: Redactar "Pruebas".

\subsection{9.2. Flujo de Trabajo de Pruebas en PUDS}

% [PENDIENTE]: Redactar "Flujo de Trabajo de Pruebas en PUDS".

\subsection{9.3. Técnicas de Pruebas}

\subsubsection{9.3.1. Caja Blanca}

% [PENDIENTE]: Casos de uso a probar (propuestos: CU-02, CU-03 y CU-05; ajustar).

\paragraph{9.3.1.1. Caso de uso: Enviar postulación a Beca Comedor (CU-02).}

% [PENDIENTE]: Redactar "Caso de uso: Enviar postulación a Beca Comedor (CU-02)".

\paragraph{9.3.1.2. Caso de uso: Precalificar postulación con Inteligencia Artificial (CU-03).}

% [PENDIENTE]: Redactar "Caso de uso: Precalificar postulación con Inteligencia Artificial (CU-03)".

\paragraph{9.3.1.3. Caso de uso: Auditar y adjudicar becas (CU-05).}

% [PENDIENTE]: Redactar "Caso de uso: Auditar y adjudicar becas (CU-05)".

\subsubsection{9.3.2. Caja Negra}

% [PENDIENTE]: Mismos casos de uso que en caja blanca (ajustar).

\paragraph{9.3.2.1. Caso de uso: Enviar postulación a Beca Comedor (CU-02).}

% [PENDIENTE]: Redactar "Caso de uso: Enviar postulación a Beca Comedor (CU-02)".

\paragraph{9.3.2.2. Caso de uso: Precalificar postulación con Inteligencia Artificial (CU-03).}

% [PENDIENTE]: Redactar "Caso de uso: Precalificar postulación con Inteligencia Artificial (CU-03)".

\paragraph{9.3.2.3. Caso de uso: Auditar y adjudicar becas (CU-05).}

% [PENDIENTE]: Redactar "Caso de uso: Auditar y adjudicar becas (CU-05)".

\subsubsection{9.3.3. Pruebas de Vulnerabilidad}

\paragraph{9.3.3.1. OWASP ZAP.}

% [PENDIENTE]: Redactar "OWASP ZAP".

\paragraph{9.3.3.2. Prueba de vulnerabilidad manual.}

% [PENDIENTE]: Redactar "Prueba de vulnerabilidad manual".

\subsubsection{9.3.4. Lista de Comprobación}

% [PENDIENTE]: Redactar "Lista de Comprobación".

\subsection{9.4. Representación del Tiempo de Respuesta}

% [PENDIENTE]: Contrastar con los RNF de rendimiento (carga <= 3 s; puntaje individual <= 5 s).

\subsection{9.5. Herramientas para Asistir el Proceso de Pruebas en el Software}

% [PENDIENTE]: Redactar "Herramientas para Asistir el Proceso de Pruebas en el Software".
